%\documentclass[12pt, oneside]{book}
\documentclass[12pt, twoside, openright]{book}

% ============================================
% ENCODING & FONTS
% ============================================
\usepackage{graphicx}
\usepackage{pifont}
\usepackage{amssymb}
\usepackage{enumitem}
\setlist[itemize]{nosep}
\usepackage{chngcntr}
\usepackage[colorlinks=true, linkcolor=black, urlcolor=blue]{hyperref}\usepackage[perpage,symbol*]{footmisc}
% footnote marks: symbol* gives the classical sequence (* dagger double-dagger section-mark pilcrow ...),
% perpage restarts it on every page, so the first footnote of any page is always the asterisk.
% (removed \renewcommand{\thefootnote}{} -- the empty-mark scheme is superseded; symbol* would be
% overridden by it if it stayed. symbol* also warns if a page ever exceeds the nine symbols.)

\hypersetup{
	pdftitle={Shoresh},
	pdfauthor={Alex Feldmann},   % pen name only — never the real name
	pdfsubject={Literary fiction},
	pdfkeywords={Literary fiction, memoir, interpersonal},
	pdfcreator={},                  % blanks "LaTeX with hyperref"
	pdfproducer={},                 % see caveat below
	pdflang={en}
}


\renewcommand{\footnoterule}{\hrule}
\usepackage{titlesec}

\usepackage{polyglossia}
\setmainlanguage[variant=usmax]{english}% variant=usmax loads the extended US pattern set (usenglishmax). The default (variant=us) is Knuth's original 1970s pattern list, which cannot break many common words at all -- e.g. it never finds med-i-cine. Verified with \showhyphens under XeLaTeX+polyglossia.
\setotherlanguages{german,portuguese}% hebrew is declared separately at the very end of the preamble -- see the note there (bidi load order)

\usepackage[autostyle=true]{csquotes}
\setquotestyle{american}
% (removed duplicate \setmainlanguage[variant=us]{english} -- the language must be set once, with the variant on that one call, above)

% (removed \hyphenation{med-i-cine e-nough}:
%   med-i-cine -- now found by the usmax patterns themselves;
%   e-nough -- was inert: exception entries are filtered through \lefthyphenmin=2/\righthyphenmin=3 exactly like
%   pattern results, so a break after a single letter is never produced. Verified by test. US style never breaks
%   after one letter anyway.
% If a word ever needs an exception, declare it at \begin{document} so it binds to the active patterns:
%   \AtBeginDocument{\hyphenation{some-word}}
% Note the same 2-before/3-after minima still apply to exceptions. To force a break the minima would forbid,
% use an explicit discretionary in the text instead, e.g. chron\-ic -- explicit \- bypasses the minima.)


% ============================================
% FOR RENDERING HEBREW -- NEEDS INSTALLATION OF CULMUS
% ============================================

\usepackage{fontspec} % for xelatex or lualatex

\setmainfont{Libertinus Serif}
\setsansfont{Libertinus Sans}
\setmonofont{Libertinus Mono}

\usepackage{microtype}     % spacing/kerning refinement; do not skip
\linespread{1.05}          % the extra leading Libertinus wants
\frenchspacing             % single space after periods — standard in books



% 
% (removed the one-line \newfontfamily\hebrewfont declaration that stood here. Current fontspec throws
% "Command \hebrewfont already defined" on the duplicate, and the detailed declaration below then never
% takes effect; older fontspec versions silently let the second override. Either way, one must go, and the
% detailed one is the keeper: Frank Ruehl CLM's faces are named Medium/Bold/MediumOblique/BoldOblique --
% no "Regular" or "Italic" -- so the explicit mappings below pin the right faces on any system.)
\newfontfamily\hebrewfont[
Script=Hebrew,
UprightFont={Frank Ruehl CLM:style=Medium},
BoldFont={Frank Ruehl CLM:style=Bold},
ItalicFont={Frank Ruehl CLM:style=MediumOblique},
BoldItalicFont={Frank Ruehl CLM:style=BoldOblique}
]{Frank Ruehl CLM}

% NEEDED IN OTHER SYSTEMS: sudo apt install culmus

% the one that reads {\hebrewfont עזאזל}, in small discreet letters




% ============================================
% PAGE GEOMETRY (6" x 9" Amazon KDP)
% ============================================
\usepackage[
paperwidth=6in,
paperheight=9in,
top=0.75in,
bottom=0.75in,
inner=0.75in,
outer=0.5in,
headheight=0.25in,
headsep=0.3in
]{geometry}

% ============================================
% TYPOGRAPHY & SPACING
% ============================================
\usepackage{setspace}
\onehalfspacing

% Prevent orphans/widows
\clubpenalty=10000
\widowpenalty=10000

% Better paragraph spacing
\parskip=0.3\baselineskip
\parindent=0.20in

% Relief valve for the narrow 6in measure: allows a little extra interword
% stretch, but only on paragraphs TeX cannot otherwise break acceptably.
% Raise toward 3em if Overfull \hbox warnings persist in the log.
\emergencystretch=1.5em

% ============================================
% HEADERS & FOOTERS
% ============================================
\usepackage{fancyhdr}
\pagestyle{fancy}
\fancyhf{}
\cfoot{\thepage}
\renewcommand{\headrulewidth}{0pt}
\renewcommand{\footrulewidth}{0pt}

% ============================================
% CHAPTER & SECTION FORMATTING
% ============================================
\titleformat{\chapter}
{\normalfont\LARGE\bfseries}
{\thechapter.}
{0.5em}
{}
[\vspace{0.5em}]

\titleformat{\section}
{\normalfont\Large\bfseries}
{\thesection}
{0.5em}
{}

\titlespacing*{\chapter}{0pt}{1.5em}{1.5em}
\titlespacing*{\section}{0pt}{1em}{0.5em}

% (removed the first tocloft block that stood here -- the fuller duplicate below, under "TABLE OF
% CONTENTS FORMATTING", is kept. A repeated \usepackage of the same package is skipped by LaTeX as
% long as the options match, but becomes an "Option clash" error the moment an option is added to
% either copy.)

% ============================================
% DOCUMENT INFO
% ============================================
\title{Shoresh}
\author{Alex Feldmann}
\date{}

% ============================================
% TABLE OF CONTENTS FORMATTING
% ============================================
\usepackage{tocloft}
\setcounter{tocdepth}{0}
\renewcommand{\cftchappagefont}{}  % Remove page numbers
% NOTE (unchanged): the line above does NOT remove the TOC page numbers -- it only strips their bold
% styling (tocloft's default \cftchappagefont is \bfseries); verified by test compile. If no page numbers
% in the TOC is what you actually want, enable the next line instead and delete the one above:
% \cftpagenumbersoff{chapter}

% Hebrew is declared last on purpose. An RTL language makes polyglossia load the bidi package
% immediately, and bidi insists on being the last package loaded: it raises errors at the start
% of the document for every listed package loaded after it ("Oops! you have loaded package
% geometry after bidi package" -- and the same again for fancyhdr). Declaring hebrew here,
% after every other package, loads bidi last and satisfies its check.
\setotherlanguage{hebrew}

\begin{document}


	% ============================================
	% FRONT MATTER
	% ============================================
	\frontmatter
	
	%---TITLE PAGE---
	\thispagestyle{empty}
	\begin{center}
		\vspace*{1.5in}
		
		{\Huge\bfseries Shoresh}
		
		\vspace{1in}
		
		{\Large Alex Feldmann}
		
		\vspace{3in}
		
	\end{center}
	
	\clearpage
	
	%---COPYRIGHT PAGE---
	\thispagestyle{empty}
	\vspace*{7in}
	
	\begin{center}
		\footnotesize
		
		Copyright \copyright\ 2026 Alex Feldmann
		
		All rights reserved. No part of this book may be reproduced, stored, or transmitted in any form or by any means without the prior written permission of the author, except for brief quotations in critical articles or reviews.

		\vspace{1em}

		First edition, 2026

		ISBN: [to be assigned]

		\vspace{1em}

		The events described are true fiction. Any similarity to actual persons, living or dead, or to actual events is purely coincidental.

		\vspace{1em}

		Set in Linux Libertine, 12/18.

		Composed in Xe\LaTeX.

		Cover photograph and design by the author.
				
	\end{center}
	
	\clearpage
	
	%---TABLE OF CONTENTS---
	\tableofcontents
	
\clearpage
	
	% ============================================
	% MAIN MATTER
	% ============================================
	\mainmatter
	
	%---PROLOGUE---
	\chapter*{Prologue}\label{prologue}
\addcontentsline{toc}{chapter}{Prologue}

%\vspace{2em}

From Hebrew:

Shoresh /  {\large שורש }

\begin{quote}
	$\bullet$ Literally: root---the part of a plant, generally underground, that anchors and supports the plant body, absorbs and stores water and nutrients, and in some plants is able to perform vegetative reproduction

	$\bullet$ by extension, root: a source or basis

	$\bullet$ by extension, in linguistics: root of a Semitic word

	$\bullet$ by extension, in mathematics: root as in square root
\end{quote}

\bigskip

Ilan /  {\large אילן}

\begin{quote}
	$\bullet$ Noun: tree.
	
	$\bullet$ Proper noun: a male given name.
	
\end{quote}

\bigskip

Nitzanah /  {\large ניצנה}

\begin{quote}
	$\bullet$ Noun: a blossom; flower bud.
	
	$\bullet$ Proper noun: a female given name.
	
\end{quote}

\bigskip
	

	
	% ============================================
	% CHAPTERS IN READING ORDER
	% ============================================
	\chapter*{Start}\label{start}
\addcontentsline{toc}{chapter}{Start}

\enquote{\noindent\textsc{{How many kids would you have had?}} I mean, if it were not for\dots{}}

Her voice slowly fades out, her face shying away the same way she does every time the topic of her mother is approached. So many times, so many years, and it still pains me to see how she gets the guilty blow-back from herself every time she does it, and yet she cannot avoid repeating it.

I say, \enquote{Hey Ladybug, don't get blue---it does not hurt like before.} It's a lie, it still hurts. And she senses it somehow.

I reach out for her, bringing her to a hug as she shyly and yet dramatically trudges into the living room from the apartment's entrance. I sniff her hair for a while and notice: smoke. I ask, \enquote{Were you smoking?}

She takes some distance, opens her mouth demonstratively and exhales in my direction: negative for smoke. 

\enquote{It was Justyna,} she says. \enquote{We gathered in the park, she's got the confirmation of a position. She will start right after finishing training, in some months. She was musing that, with her and Thomas having confirmed positions, they can start to make plans for a baby.}

\enquote{So, \emph{this} is where this kid-number-story comes from,} I comment. \enquote{Now tell me: how will Justyna have a baby and smoke? Or\dots{}}

\enquote{\dots~she will have to stop when she decides to get pregnant, so she wants to smoke now while she can,} she complements my line of thought. She is good at doing that---just like her mother. I usually hate when people try to predict what others are about to say, but only because most of the time people get it wrong. 

Nitzanah never does it. 

I invite her, with gestures, to sit down. She does it, drops her backpack near the couch. She still dresses like a kid, even though she's a woman old enough to muse about having kids in earnest. I feel glad that I enjoyed her years as a kid as fully as I could. 

I decide that it starts now. 

\enquote{Nitzy\dots{}} She understands by the undertone that her attention is required.

\enquote{Yes?}

\enquote{Did Justyna comment on pre-natal precautions, like genetic screening?} I ask.

\enquote{No\dots~is that not only recommended for persons at risk of some problem?}

I answer, \enquote{Mostly, yes---and I do not know if she or her partner---Thomas?---are at risk\dots~but perhaps it is important that you know about us.}

Her whole posture and attitude changes, now she's straightened up and looking at me very earnestly. She's waiting for the continuation. I abide.

\enquote{For one, I do not believe I told you before, but I am a thalassemia carrier. It is a genetic condition that can cause anemia of diverse degrees, depending on how many genes are affected. Now, you do not have anemia, and probably do not have any gene affected, but there are several conditions that may run silent and still be a problem for the children. And there are things about your mother's side you do not know, but should, before you have children.}

She looks at me for a while. After some reflection that I recognize from oh so many sweet moments, she declares, \enquote{You are preparing me to get away from you. To be alone. Stop.}

I laugh out loud and reply \enquote{Of course I am not! Do not be so dramatic. Just like your mother, aren't you? Just look at you: a fully grown woman. If Justyna can start having kids anytime, so do you. And if you do, why not check things beforehand?}

She smiles at me in a warm way. It breaks my heart not being capable of telling the truth as it is and hiding behind a lab result.
	\chapter*{Bloodwork}\label{bloodwork}
\addcontentsline{toc}{chapter}{Bloodwork}

Funny, actually funny. A different kind of clinic, not like the routine visits to the immunologist. The nurses look different, laugh differently, behave differently--- they just \emph{are} different.

Dad waits reading at the reception, the whole procedure is rather simple. I imagined that they would take half of my blood away, like, a big bottle or something, but it does not look like much. I guess that even the immunologist claimed more blood from me than what these guys are doing.

I feel positive the whole time and it is just when I meet dad again, when he raises his eyes from the magazine he's reading, that I notice something. Again, like he is saying goodbye to me. 

I try to find a way to him, but I cannot. I want to believe that none of us is responsible, but I have no foundation for that.
	\chapter*{Result}\label{result}
\addcontentsline{toc}{chapter}{Result}

I try to behave normally but I feel that she notices that something is not quite alright. I try to divert her attention, I falsely ask, \enquote{Hey, Ladybug, why are you so nervous? Remember, this is more about possible kids than it is about you\dots{}} and I instantly regret it and want to take it back.

She looks at me with an inquiring look and gives me no answer---she is not nervous at all. I am.

I will regret these last moments of innocence forever, because I am tarnishing them. 

\bigskip\noindent\textsc{The news hits like a stone falling on stamped earth.} The physician looks at me with the knowing look that I was expecting. Nitzanah silences at my side the way I expected, for long moments. The physician is friendly enough to give us time. 

I do not dare to look at her. Like the coward that I am, I hide behind the armrests of my chair. Minutes, hours, days go by---the physician is understanding enough to have postponed all his other appointments.

She is clinging to me and I weep. A nurse comes in, the physician orders something. We are both brought to another room. I take the medication the physician ordered---she does it too. It will help, the nurse says. 

\bigskip\noindent\textsc{An eternity goes by.} The clinic's personnel orders a taxi, we are brought home. The driver follows us until the apartment's entrance and makes sure we get inside. A patient's taxi driver.

Now the most dreaded moment: to be alone with her.
	\chapter*{Talk}\label{talk}
\addcontentsline{toc}{chapter}{Talk}

He wants to hide but I won't let him. I block the narrow corridor that leads to the apartment's kitchen and rooms. I demand answers, I want them now, but he won't do it now. He says he will explain everything later.

He confirms that he knew it, all along. I spit out the truth, that he never told me, never hinted at anything. He replies that there would never be an ideal moment and my adulthood presented the best opportunity. I feel like I was a puppet being pulled by strings this whole time. I know that parents manage their kids and don't always tell everything, but this is too much.

He asks for time, to rest, to bring \enquote{order to his thoughts.} What about my thoughts, my feelings? I am the victim here, I remind him. I see that it hurts and I feel a jolt of joy but also a sting of sadness. I give him the time he wants.

I hear him call in sick at work. He promises to have the whole following day free for talk.

\bigskip\noindent\textsc{The news hits like a stone falling on stamped earth.} The physician looks at me with the knowing look that I was expecting. Nitzanah silences at my side the way I expected, for long moments. The physician is friendly enough to give us time. 

I do not dare to look at her. Like the coward that I am, I hide behind the armrests of my chair. Minutes, hours, days go by---the physician is understanding enough to have postponed all his other appointments.

She is clinging to me and I weep. A nurse comes in, the physician orders something. We are both brought to another room. I take the medication the physician ordered---she does it too. It will help, the nurse says. 

\bigskip\noindent\textsc{An eternity goes by.} I sleep through the afternoon, she apparently too. One might assume that a genetics clinic has preparedness for stressful revelations of all kind. The clinic's personnel orders a taxi, we are brought home. The driver follows us until the apartment's entrance and makes sure we get inside. A patient's taxi driver.

Now the most dreaded moment: to be alone with her.
	\chapter*{Explanations}\label{explanations}
\addcontentsline{toc}{chapter}{Explanations}

I understand that she needs time. I have been bracing for this moment since I made the decision, nearly two decades ago. I always knew this moment would come. And I never imagined how unprepared I would feel when it finally came.

I feel physical pain in my body, a pain that comes from seeing her suffering, and from knowing that I am the reason of her suffering.

I approach her from behind, she is giving me her back, watching through the window and I know why. I touch her shoulders lightly, just to deliver her a tactile confirmation that I want to be with her. 

I tell her, \enquote{Hey, there is a lot of things that I can explain to you. First, my reasons, of course. But also about the technical side. I would prefer that you get it from me than from some internet source.}

She stays silent for some moments. Then a reply, \enquote{Okay. Just give me some moments.}

And she goes to gather a paper notebook and a pencil. I take my usual seat at our small wooden kitchen table while she scribes vigorously one bullet point after the other. I feel proud of her. 

I look at her and I wonder if the fact that she looks so little like her mother is why I love her so much. The straightness of the hair does not match, the facial structure is different, probably a consequence of repeated rhinitis and sinusitis during childhood, even the height and body build differ. So many things can influence how bodies develop, what to say about minds?

She wakes me up from my daydream, asks, \enquote{So, what is my legal situation? In Germany?}

Straight to the point---that's my girl. I tell her what I always knew and kept up-to-date, \enquote{German law consider all human beings worth of human dignity, independently of how they come into existence. The process by which you came to life is not allowed in Germany, but you had no part in it and cannot be blamed. If at all, I would be the criminal, if it were done here\dots~but it was not, it was in China, and there, it is legal.}

\enquote{So, none of us is going to prison?}

I reassure her, \enquote{No, neither you, nor me. And the genetics clinic is now allowed to share the information with anyone else without your permission.}

She looks genuinely relieved. She asks, \enquote{Are you sure about that?}

I answer, \enquote{I was sure even before I started everything. I tend to think long-term, Ladybug.}

She gives me the first smile since the revelation, since the result.

\enquote{Okay, next question. But before the question: I understand that you needed to wait until I was\dots~old enough to understand the situation. But then the question: why to mention her as my \enquote{mother} when she is actually something different? I do not have a word for it\dots{}}

I reply, \enquote{\emph{Shoresh}---I suggest that we use that. It is Hebrew for \emph{root.} I guess that it is the best description, in several instances.}

She looks down for several moments. When she looks back at me, there is understanding in her eyes. 

She tells me, \enquote{Yes, I have a father, but I have no mother---I have \emph{Shoresh}. I can live with that. I must live with that. But I want to know who she was. Tell me.}

So I do as requested. I do it calmly, because time is not a problem now. Before the revelation, I feared and dreaded this moment, but now that the worst part is over, I want to savor these moments with her. I want to do it the right way.

I brew us a large pot of tea, the one that we use when there are visitors. And I grab a whole package of oat cookies along with a bar of dark chocolate and some dry fruits.

I assemble all things together on the small rustic wooden center table at our living room and call her to join me. I choose a playlist that, in retrospect, seems that was built exactly for this moment. 

	\chapter*{Acquaintance}\label{acquaintance}
\addcontentsline{toc}{chapter}{Acquaintance}

He does it well. I can see that he's been preparing for this. He brings the cookies that I once liked but have for a while overgrown. He wants to make me comfortable. 

I cannot keep hating him like I was doing just until now.

He starts from the very obvious: how they met each other. It does not interest me. I do not think that people are themselves when they are meeting someone new. I am not myself when meeting someone new.

I try not to interrupt, it is obvious that he is heading toward what I whish. 

Now: her first personal wish. 

Why did she want to go there?

Because it was the city of her family. 

This does not make sense, Jews were not Germans citizens.

At some point, they had to become, Nitzy, it was enforced law.

So everything was un-enforced decades later?

Ironies of history, but, yes\dots

And what happened there?

Not much, we visited the place where her family lived, and talked about what it meant.

Silence; tears; I don't know why, I do not know what to make of that, because she---\emph{Shoresh}---had interests that I cannot understand, and that he cannot communicate.

How many more, of relatives, I mean?

Not many, there is a branch of the family, but it is not close.

Have you had contact with them?

No, never, they emigrated to the Americas.

Why not?

Hey, this was her family, I do not know everything.

Go on, tell me more.

Wait, don't you want to visit the location?

Do not distract me, we have only today for information, but forever for places. 

He seems struck. I regret saying that, it sounds like a type of goodbye.

She's chosen your name.

What?

Yes, we talked often about that, which names to choose in case\dots

So you were\dots~planning on having a child?

Yes, we even got genetically tested.

I have to think for a while.

So the whole talk about thalassemia\dots

\dots~was already known. Yes, both me, and her, had it, and so do you.

This makes me dizzy, I need to ask now.

So, it might be that even if she had not died, you would have opted to have a kid this way, as you yourself did alone?

He takes his time to answer.

\enquote{Yes.}
	\chapter*{Acceptance}\label{acceptance}
\addcontentsline{toc}{chapter}{Acceptance}

I leave her busy with her thoughts, as I imagined she would be after the conversation. I tell her that I am open to proceed further, but also that I will share with her some resources to orient her in the flood of information that the internet provides. It does not help that very little public facing information is in Germany, because the procedure is not allowed here and in the European Union.

But what really annoys me is the amount of information that is misleading or outright false, mostly from ads promoting clinics in other countries. It is inevitable that she will see the bad websites, sooner or later, but I would like her to read my selection first. 

I do not share all at once; I give her at first enough reading for a part of today, so we can discuss after she finishes. As part of my curation, I share scientific and public facing articles on identity dysphoria disorders, which can vary from over-identification to consciously targeted negation of any similarity with the root. 

She reads in her tablet, still in the kitchen, and does not care to sit on the couch or lie on her bed. I stay around in case she has any questions and pretend to read something in my device, but I actually just watch her. Of course she looks a lot like her mother, even if I tried to convince me otherwise, but not like I feared she could become when she was young. The most striking similarities I notice are actually in her behavior, such as now, how deeply she can concentrate on something, and that flickering in the eyes that is the same.

It pains me to know that I will never be able to talk about that with her.

\bigskip\thinspace\noindent\textsc{After she finishes the readings I have lunch ready.} We eat it in silence and I suggest that we go for a walk in our neighborhood. I do not tell why, but it is to anchor her through the familiar sights, to avoid dissociating from the life she had until now. To reassure her that she is still the same person.

It seems to work, I would like to think. Several times she stops and looks at sights that she saw thousands of times before, and I wonder if she is seeing something new or acknowledging the old. I do not ask, I fear the answer.
	\chapter*{Understanding}\label{understanding}
\addcontentsline{toc}{chapter}{Understanding}

He did well. I realize now that he has been preparing for this moment all my life. The revelation via medical examinations might not have been planned, but was an opportunity well used.

The information package was also very helpful. Planned to be explanatory enough. For a layperson. Standing in the park, now, might seem spontaneous, but I am sure he had it planned.

I feel grateful but also hateful. Grateful wins. I decide to talk again. I point to a bench under a tree. 

Correct me if I am wrong.

Right.

My\dots~mind? My mind is my own because brains develop differently.

Right. Remember the comparison with twins, who are undoubtedly different persons. 

So it's like she is my sister, not my mother.

Right. If one of the twin embryos were frozen for years and only later implanted, it would be the same.

Nothing in the information package is as simple as that. Why?

Because people make it more complicated than it is, because of social, psychological, religious aspects. No wonder many countries still prohibit it. So many people still think that it is about making \enquote{copies,} or even worse, \enquote{backups} of oneself. But as someone who knew you both I can assure you: you both have very distinctive personalities.

You never talked much about her. Was it to avoid influencing me?

Not exactly. I did not do it at first because it hurt\dots~and later on because you in my life was all I needed. I did not forget her, but there was no longer a painful void.

I guess that it is enough for now. Not forever, just for now. May I ask again anytime?

Of course! Anytime, anything.

I lean on his shoulder for a while. I can feel his heartbeat. It was faster right after the talk, then got slower. He takes my hand, presses it lightly. I try to quantify how much more he had to invest than usual parents. Personal effort, money, sacrifices in his career. I cannot quantify.

I feel overwhelmed by the thought. My arm thrown over his shoulders bringing him closer tells him that.  
	\input{chapters/Adaptation}
	\chapter*{Probing}\label{probing}
\addcontentsline{toc}{chapter}{Probing}

She brings no more questions that first day. After a while in silence in the park we walk home and she spends the rest of the day, and evening, on her bed. Sometimes I catch her just watching the ceiling, sometimes she let music play. When we have dinner there is just functional dialogue about the food, about news. But it is clear to me that her head is busy, it is only that she still does not have the words for it.

On the following day she has the part-time junior job with which she occupies herself until classes at the engineering school started, thereafter she will meet friends and spend time in the old city. Before leaving she tells me that, and holds my gaze for a while, watching my reaction. She is silently asking me if it was okay to talk to her friends about it.

I asked in return, \enquote{They are your friends, you know them; how do you think they would react?}

\enquote{I don't know, we never talked about it\dots{}}

\enquote{Then you can probe it, mention something that you read from that package I gave you. See how they react. But most importantly, think about \enquote{why} you want to bring up the topic with them. Do you want it to make some difference? Or do you want to be sure that it makes no difference?}

She left as thoughtful as she was before the talk, her eyes lost in the infinity that lies under the floor in front of each one of us.
	\chapter*{Care}\label{care}
\addcontentsline{toc}{chapter}{Care}

I did not expect it so soon. I was aware that, at some point, some crisis might erupt, but this is just the third day.

She came back from the time with her pals, greeted me in a perfunctory way, went to her room and is now lying on her bed, on her right side, looking at the emptiness of the wall in front of her. No crying, no movement. I take a look regularly, probably each half an hour, but I see no change.

Around twenty I bring her tea, the type she has preferred lately, also some nuts and dried fruits. She likes that, and it is energy food that does not require much on her side.  

When it is half past twenty-two I come into her room and sit by her side. The whole time she has not moved. 

I look at her face, which shares so much of her mother, and yet has so much of her own. I seem to see also something of mine there, but I know that it is probably an illusion, or, at best, some mannerism which she learned from me. 

Looking at her I ask of myself, was it fair? To bring her into this world in this way? Was I being so selfish as to not imagine that one day she would suffer?

I say to myself that I would die if that would spare her the suffering she is living through now, but that is a cheap saying, because it is impossible to prove it true. Worse than that, it would mean abandoning what I did.

There is only one way forward, and it is together with her. 

She seems to listen to my thoughts, she turns her face lightly to me. My face asks something. Her face expression tells me \enquote{Not yet.} 
	\chapter*{Crisis}\label{crisis}
\addcontentsline{toc}{chapter}{Crisis}

I come back, lie in bed. Stay there. He comes, goes, brings me food, tea.

Next morning he comes again. I turn to him, I say, I always knew something was wrong.

It is a lie 

I always knew you were lying.

It is a lie.

I am not who I thought I am. 

It is true. He gives me no answer because it is not possible. I hear him calling my work, his work, of.

I say, no. I need this. I give up.

He comes closer, caresses the back of my neck, like aways. He whispers, this changes nothing about who you are. I love you. You, not genes or cells.

I shrink down to the smallest I could, until I fit into his hand. I slow time to the slowest I can, until all thoughts almost stop. I want to stop time for forever but I fail. I want to become so small that I vanish but I fail.

I want to fail to keep existing. But even at that want, I fail.

\bigskip\thinspace\textsc{New lights wake me up.} My eye sees a new world. New sounds come through the window. A new pair of breathe in and out. A new heartbeat gives me rhythms. 

New arms were holding me, still asleep. He is not new; I hold him closer, afraid that he can get too small and slow to love me. He cannot, he is my father. 

I am there with him until the universe ends. I see all stars shutting down. I am alone with him and that feels correct. I wish this could last forever. It does not. 

\bigskip\thinspace\textsc{He brings me food several times}. I do not refuse it but I do not see the reason. He leads me to the toilet several times but I do not see the reason. He brings me water and milk and tea. He did it several times, but I do not see the reason. I opened my eyes when he did all of that and I see it all but I see no reason. All several times happened several times, then this also several times and I see no reason.

He holds me me. I shrink down to be smaller as a human heart. I do not expect a reason.

\bigskip\thinspace\textsc{I start to move but it is no help.} I need to see where to go, but I see no light. I bump onto things, then I give up. One time I grasped something and identified it: Father. I wished that he would love me again, help me. I wished that he could see my wish. 

But his eye is not like mine. I wish he had my eye but I regret wishing this because my eye is not really mine and I do not want him to feel what this is. 

\bigskip\thinspace\textsc{A sound that matches the rhythms of my heart.} I see, smell it. I try to grab it and fail. I try some more and fail. I am trying the wrong way: I try with my hands, then I find it. I open it up and it told me something: Nitzy. 

Nitzy, dear, what did you say?

Father, help me.

Yes, yes, stay with me.

Please, do not go away.

Yes, I am here; look here, look at me. Nitzy, I am here, never away.

Where are we?

We are home. Our home. Stay here, let me hold you.

Why are you here?

Because I love you, remember this. Keep this.

I do not want to go away.

You never will, Nitzy. Can I show you? Let me show you.

I felt his arms around me. I have tears now. I feel other tears also. I recognized them: they smelled like Father's.


\bigskip\thinspace\textsc{I feel so weak,} like there is no body down there. He's being so being so careful as if he was carrying a baby. I feel like a baby.

I feel sorry for him. I wish I would feel worth enough to feel sorry for myself.

Feeling sorry for him means motivation: I want to honor his effort. 

I bring him close, when and if he let me do it. He caresses me, as a kind of payment. I feared caressing him, because it would not be enough.

At first I felt it not. I saw him doing, only.

I got used and then I felt them. They ended in tears, but I did not regret it. It was good tears.

At some point I forget the fear and touch him without touching. He senses even though I try to be as light as I could. He smiles.

%\begin{center} \rule{0.3\textwidth}{0.4pt} \end{center}
%\begin{center}
	\vspace{1em}
	\rule{0.15\textwidth}{0.4pt}\hfill$\ast$\hfill\rule{0.15\textwidth}{0.4pt}
	\vspace{1em}
\end{center}
%\begin{center}
	\vspace{1em}
	\rule{0.2\textwidth}{0.4pt}\quad{\scshape §}\quad\rule{0.2\textwidth}{0.4pt}
	\vspace{1em}
\end{center}
\begin{center}
	\vspace{1em}
	$\diamond\ \diamond\ \diamond$
	\vspace{1em}
\end{center}
%\begin{center}
	\vspace{1em}
	\ding{167}\quad\ding{166}\quad\ding{167}
	\vspace{1em}
\end{center}

	\vspace{1em}

%\begin{center}
	\vspace{1em}
	{\large \textasteriskcentered\quad\textasteriskcentered\quad\textasteriskcentered}
	\vspace{1em}
\end{center}


\emph{Frühling}\thinspace\footnote{\emph{German}: spring.} in Germany brings light to everyone: the very name sounds like it. I hear it through the window before I open my eyes.

When do I open my eyes, Mal is by my side, looking at me. He says, \enquote{Good morning.}

I smile a \enquote{Good morning.}

\enquote{How are we today?}

\enquote{Better than yesterday, worse than tomorrow.}

\enquote{I love you.}

I say, \enquote{I know that.}

\enquote{And?}

\enquote{I love you too.}

Mal points out, \enquote{If you stop saying it, I will forget it. Do you want that?}

I smile a brighter smile, \enquote{No, I don't.}

\enquote{Keep that in mind and practice it. Meanwhile, I will do breakfast. Do not try to run away; if you do, tell me before.}

\enquote{Yes, can do.}

After coming back from sleep, my thoughts still need some time to catch up with the body; I can observe it backwards. Mal also.

As I watch the ceiling, I realize, like I have been doing in the last mornings, how much of a toll this must have been taking on Mal.

After we have breakfast, he will go to work and check on me during the day. He is lovely.

I wish I could assure him that I need less intensive monitoring and care, but I realize I am not capable of changing this. I will have to show it through the facts and results.

So I send him pictures of me:

\begin{itemize}
	\item smiling while feeding the cat of the neighbor who is away (this is
	not a lie)
	\item enjoying doing workout at home (this is a lie)
	\item random things like varied city landscapes, buildings, etc (this is not
	a lie)
	\item sunbathing and reading books in sunny parks (they are not a lie)
\end{itemize}

I will not let anyone judge me on lying: if anyone knew him as I do, they would do the same.

I do not try to play being fully okay, and he would notice if I did. But what is amiss, I still cannot see what it is. But I am sure I will; I might not do it alone---but I am not.

%\begin{center} \rule{0.3\textwidth}{0.4pt} \end{center}
%\begin{center}
	\vspace{1em}
	\rule{0.15\textwidth}{0.4pt}\hfill$\ast$\hfill\rule{0.15\textwidth}{0.4pt}
	\vspace{1em}
\end{center}
%\begin{center}
	\vspace{1em}
	\rule{0.2\textwidth}{0.4pt}\quad{\scshape §}\quad\rule{0.2\textwidth}{0.4pt}
	\vspace{1em}
\end{center}
\begin{center}
	\vspace{1em}
	$\diamond\ \diamond\ \diamond$
	\vspace{1em}
\end{center}
%\begin{center}
	\vspace{1em}
	\ding{167}\quad\ding{166}\quad\ding{167}
	\vspace{1em}
\end{center}

	\vspace{1em}

%\begin{center}
	\vspace{1em}
	{\large \textasteriskcentered\quad\textasteriskcentered\quad\textasteriskcentered}
	\vspace{1em}
\end{center}

I had him arrive earlier and worried, but it should be worth the trouble: I had prepared for him his favorite hummus recipe, with the marinated Aubergines he loves more than me, with garlic olives, all with \emph{Rotwein}\thinspace\footnote{\emph{German:} red wine.} of his favorite local winery and his hateful \emph{Vollkornbrot}; I am recovered enough to survive it.

For the evening, I chose a different outfit, and was gladly finishing last touches to the food when he arrived. Things from old times, when I was younger and experimented with dresses and things for fun, one of the few that I selected to save from the \enquote{old life.}

Mal entered our kitchen and froze. It took me milliseconds to notice something was wrong. Then it took me even less milliseconds to notice that what I was wearing were the same yellow cotton dress and a silver butterfly pendant that I chose for her to meet him at his second arrival at Tel-Aviv. Mal was sitting on a chair, looking at me from below with teary eyes.

I closed my eyes and clamped my ears and started crying as if ignoring things would undo them. I was falling down again---but now something was different; I felt arms around me and I could run my hands around them; I tasted tears not mine and they smelled like her. I opened my eyes and he was not visible because we were one and I can only see myself with a mirror.

Yes, on his birthday Mal chose me. Or he was just diverting attention from himself, which he always does, which ends same way.

He did not make it openly, but feel did I do. Already before this last crisis I was trying to find a way to get a trade in Germany. The official ways here expect a traceable life path, which is nowadays increasingly sensitive to scrutiny.

The way he conducted the interaction awoke something in me. He asked questions for which he knew the answers. He nodded in acknowledgment as if he was surprised. He smiled as if it meant only appreciation. It might be strange to anyone who was never there down deep that having ideas was a remarkable thing---but they are.

%\begin{center} \rule{0.3\textwidth}{0.4pt} \end{center}
%\begin{center}
	\vspace{1em}
	\rule{0.15\textwidth}{0.4pt}\hfill$\ast$\hfill\rule{0.15\textwidth}{0.4pt}
	\vspace{1em}
\end{center}
%\begin{center}
	\vspace{1em}
	\rule{0.2\textwidth}{0.4pt}\quad{\scshape §}\quad\rule{0.2\textwidth}{0.4pt}
	\vspace{1em}
\end{center}
\begin{center}
	\vspace{1em}
	$\diamond\ \diamond\ \diamond$
	\vspace{1em}
\end{center}
%\begin{center}
	\vspace{1em}
	\ding{167}\quad\ding{166}\quad\ding{167}
	\vspace{1em}
\end{center}

	\vspace{1em}

%\begin{center}
	\vspace{1em}
	{\large \textasteriskcentered\quad\textasteriskcentered\quad\textasteriskcentered}
	\vspace{1em}
\end{center}

\chapter*{Back on feet}

\enquote{It is beautiful---wonderful,} I said.

Mal smiled pleased, said: \enquote{Glad you liked it.}

I retorted, \enquote{But it was a risky move.}

\enquote{Well worth seeing your face now.}

It was a piece of jewelry that I had never seen before---or if I had seen, unconsciously suppressed: a very realistic prosthetic eye, nearly identical to my natural one, mounted into a silver socket, sculpted eyelashes, making a pendant, on a simple silver chain with round links. I wonder how he did it---I do not remember him taking any photos of me, he knows that I do not like being photographed.

I do not know if he intuited what I then realized, but I noticed at that moment that even after all those years of practice I still had some wariness, self-consciousness about my eyes. Wearing it, looking at the mirror, felt like this was what was missing.

Oh, how lucky I am to have him.

%	Epilogue included per decision in progress. The shim below guarantees a fully blank
%	leaf between "I am Helena." and the letter in BOTH page parities (openright alone
%	only separates them when Lighthouse happens to end on a recto). To re-exclude the
%	Epilogue, comment out the two lines below.
	\cleardoublepage\thispagestyle{empty}\null
%	\input{chapters/Epilogue}

	\backmatter
	
\end{document}
